\chapter{Theoretische Grundlagen}

\section{Bildverarbeitung}

Die Vorlesungsunterlagen zur industriellen Bildaufnahme und Bildverarbeitung behandeln für
dieses Projekt besonders relevante Themen: Bildaufnahme, Bildqualität,
geometrische Transformationen, Kamerakalibrierung und 3D-Rekonstruktion. Genau diese
Bausteine bilden den klassischen Bildverarbeitungsteil des Projekts.

\subsection{Kamera, Kalibrierung und Rektifizierung}

Eine reale Kamera bildet die Welt nicht ideal pinhole-artig ab. Brennweite, Hauptpunkt und
Verzeichnung müssen deshalb modelliert werden. Diese Parameter werden als intrinsische
Parameter bezeichnet. Bei einem Stereo-System kommt zusätzlich die extrinsische Beziehung
zwischen linker und rechter Kamera hinzu, also Rotation und Translation zwischen den beiden
Kamerakoordinatensystemen.

Im Projekt wird ein ChArUco-Board verwendet. ChArUco kombiniert ArUco-Marker mit
Schachbrett-Ecken. Dadurch können Kalibrierpunkte auch dann eindeutig erkannt werden, wenn
nicht das gesamte Board sichtbar ist. Das eingesetzte Board besitzt $11 \times 8$ Quadrate,
eine Quadratgröße von $10\,\mathrm{mm}$ und eine Markergröße von $7\,\mathrm{mm}$. Da diese
Geometrie in Millimetern angegeben ist, werden auch die rekonstruierten 3D-Punkte in
Millimetern interpretiert.

Nach der Kalibrierung werden Stereo-Bildpaare rektifiziert. Die Rektifizierung transformiert
beide Bilder so, dass korrespondierende Bildpunkte auf derselben Bildzeile liegen. Dadurch
wird aus einer zweidimensionalen Suche nach passenden Bildpunkten eine eindimensionale Suche
entlang horizontaler Epipolarlinien.

\subsection{Disparität und 3D-Reprojektion}

Die Disparität beschreibt die horizontale Verschiebung eines Objektpunktes zwischen linkem
und rechtem Bild. Nahe Objekte erzeugen größere Disparitäten als weit entfernte Objekte. Im
Projekt wird Semi-Global Block Matching eingesetzt. Das Verfahren vergleicht lokale
Bildbereiche, berücksichtigt aber zusätzlich Glattheitsannahmen, damit die berechnete
Disparitätskarte weniger sprunghaft wird.

Gültige Disparitätswerte werden anschließend mit der von OpenCV erzeugten
Reprojektionsmatrix $Q$ in 3D-Punkte umgerechnet. Die Farbinformation wird aus dem
rektifizierten linken Kamerabild übernommen. Ergebnis ist eine farbige Punktwolke, die als
PLY-Datei gespeichert und mit Open3D, MeshLab oder CloudCompare betrachtet werden kann.

\subsection{Bildqualität und Artefakte}

Stereo-Matching ist empfindlich gegenüber schwacher Textur, Überbelichtung,
Bewegungsunschärfe, spiegelnden Oberflächen und ungeeigneten Disparitätsgrenzen. Deshalb
enthält der Workflow Kontrastverbesserung durch CLAHE, eine automatische Warnung bei
Disparitätssättigung und nachgelagerte Filter für Tiefenbereich, Disparitätsbereich und
zusammenhängende Komponenten.

\section{Machine Learning}

Die Vorlesungsunterlagen zu Machine Learning in der industriellen Bildverarbeitung behandeln
unter anderem überwachte Lernverfahren, Datenmanagement, Convolutional Neural Networks,
Transfer Learning, Datenaufbereitung und typische Herausforderungen. Für dieses Projekt sind
vor allem CNNs, Transfer Learning, Datenlabeling und Augmentation relevant.

\subsection{CNNs und Objektdetektion}

Convolutional Neural Networks sind für Bilddaten geeignet, weil Faltungsfilter lokale Muster
wie Kanten, Ecken, Texturen und Objektteile über das gesamte Bild wiederverwenden. Durch
Parameter Sharing und lokale Rezeptivfelder können CNNs Bildmerkmale effizient lernen. Bei
der Objektdetektion wird nicht nur eine Klasse für das gesamte Bild vorhergesagt, sondern es
werden zusätzlich Bounding Boxes für einzelne Objekte bestimmt.

YOLO-Modelle formulieren Objektdetektion als schnellen Vorwärtsdurchlauf eines neuronalen
Netzes. Die Ausgabe besteht aus Klasse, Konfidenz und Bounding Box. Im Projekt wird YOLO als
semantisches Modul eingesetzt: Es erkennt Schreibgeräte in den Rohbildern, während die
Stereo-Pipeline die geometrische 3D-Information erzeugt.

\subsection{Transfer Learning und Datenaufbereitung}

Das Training moderner Detektionsnetze von Grund auf benötigt große Datenmengen und viel
Rechenleistung. Deshalb wurde ein vortrainiertes YOLO11m-Modell verwendet und auf den eigenen
Stiftdatensatz angepasst. Das entspricht dem Transfer-Learning-Gedanken: allgemeine
Bildmerkmale werden aus einem vortrainierten Modell übernommen, während das Modell auf die
projektspezifischen Klassen feinangepasst wird.

Für überwachte Detektion ist die Labelqualität entscheidend. Die Projektbilder wurden in
Bounding-Box-Form annotiert und in Trainings-, Validierungs- und Testdaten aufgeteilt.
Augmentationen wie Skalierung, Translation, horizontales Spiegeln, Farbanpassungen, Mosaic
und Random Erasing erhöhen die Robustheit gegenüber realen Aufnahmevariationen.

\section{Trennung der Projektteile}

Die Bildverarbeitung arbeitet deterministisch mit expliziten Kameramodellen und
geometrischen Annahmen. Ihre Ergebnisse sind metrisch interpretierbar, aber empfindlich
gegen Kalibrier- und Matchingfehler. Das Machine-Learning-Modul arbeitet datengetrieben. Es
liefert semantische Entscheidungen, deren Qualität von Datensatz, Labeling und
Generalisierung abhängt. Im Projekt sind beide Module getrennt lauffähig und werden erst auf
Ebene der Anwendung zusammengeführt.
